% ==============================================================================
% ПАРТИЯ B03-018: DISS-005-B0755 .. DISS-005-B0812
% Глава 2: Введение, КР и ТТ, формулы тока КЗ (2.1, 2.2)
% Замечания: DISS-005-C0074 .. DISS-005-C0076 (3 шт.)
% Объекты: 33 шт. (формулы 2.1, 2.2, 2.3, 2.4, рисунки 2.1, 2.2, 2.3, 2.4, 2.5)
% ==============================================================================

% DISS-005-B0756
\section{Введение}\label{sec:ch2_intro}

% DISS-005-B0757
Надежность электроснабжения\commentnote{DISS-005-C0075}{Алексей Евдаков [2]}{Накидано нейронкой, но мне нравится.} промышленных предприятий и других ответственных потребителей является ключевым фактором их бесперебойной работы и экономической эффективности. Системы быстродействующего автоматического ввода резерва (БАВР) играют в этом важнейшую роль, обеспечивая оперативное переключение на резервный источник питания при возникновении нарушений на основном. Эффективность функционирования БАВР напрямую зависит как от точности и быстродействия измерительных органов, так и от корректности логики восстановления нормального режима электроснабжения после устранения аварийной ситуации.\par\vspace{0.2cm}

% DISS-005-B0758
Традиционные подходы к построению токовых измерительных цепей на основе электромагнитных трансформаторов тока (ТТ) обладают рядом известных ограничений, таких как насыщение сердечника при больших токах короткого замыкания и значительное влияние апериодической составляющей, что может приводить к замедлению или неверной работе защит. В то же время, классические алгоритмы восстановления нормального режима (ВНР) не всегда способны адекватно реагировать на сложные сценарии в энергосистеме, такие как несинхронная работа источников, наличие однофазных замыканий на землю или одновременное восстановление питания на множестве присоединений, что может привести к повторным авариям или недопустимым перегрузкам оборудования.\par\vspace{0.2cm}

% DISS-005-B0759
В данной главе рассматриваются два важных аспекта модернизации устройств БАВР, направленные на повышение их общей эффективности и надежности.\par\vspace{0.2cm}

% DISS-005-B0760
Первая часть главы (2.2) посвящена исследованию применения катушек Роговского (КР) в качестве первичных преобразователей тока для устройств реле направления мощности (РНМ), являющихся важным компонентом БАВР. Анализируются преимущества КР по сравнению с традиционными ТТ, исследуется их поведение в переходных режимах, влияние различных факторов (апериодической составляющей, высших гармоник, параметров дискретизации и оконных функций БПФ) на точность измерений, а также рассматриваются практические аспекты их применения.\par\vspace{0.2cm}

% DISS-005-B0761
Вторая часть главы (2.3) посвящена совершенствованию алгоритмов функционирования восстановления нормального режима. Анализируются существующие проблемы и ограничения стандартных алгоритмов ВНР и предлагаются модернизации, направленные на предотвращение некорректных переключений при однофазных замыканиях на землю, значительных расхождениях по углу или частоте между источниками, а также для снижения рисков при одновременном срабатывании нескольких устройств ВНР.\par\vspace{0.2cm}

% DISS-005-B0762
Целью настоящей главы является демонстрация возможностей повышения надежности и быстродействия БАВР за счет применения современных измерительных технологий и усовершенствованных алгоритмов управления, что в конечном итоге способствует повышению стабильности электроснабжения.\par\vspace{0.2cm}

% DISS-005-B0763

% DISS-005-B0764
\section{Функционирования РНМ с катушками Роговского\commentnote{DISS-005-C0076}{Евдаков Алексей}{Взято из статей про КР}}\label{sec:ch2_rnm_rogowski}

% DISS-005-B0765
Применение катушек Роговского в устройствах БАВР в России началось только в 2023 году. Ранее не было каких-либо системных исследований о том, как применение данных датчиков будет влиять на функционирование БАВР и как сигналы, полученные с КР будут влиять на поведение органов реле направления мощности в переходных моментах, применяемых для определения место повреждения.\par\vspace{0.2cm}

% DISS-005-B0766
Данные изменения токовых цепей на данный момент применяются в следующих вариантах:\par\vspace{0.2cm}

% DISS-005-B0767
При модернизации существующих подстанций 6~$\div$~35~кВ с целью удешевления и упрощения монтажных работ за счёт отсутствия необходимости в подключении токовых цепей (установка дополнительных ТТ, прокладка медных проводов и т.д.).\par\vspace{0.2cm}

% DISS-005-B0768
При установки современных терминалов БАВР, предназначенных для работы на подстанциях 0,4~кВ. Это обусловлено полным отсутствием каких-либо измерительных датчиков тока.\par\vspace{0.2cm}

% DISS-005-B0769
При установке БАВР на современных подстанциях, выполненных на ячейках, осуществляющих полный перечень защит, и, следовательно, исключающих возможность подключения дополнительных цепей. Отметим, что в этом случае, дополнительно был разработан промежуточный преобразователь аналоговых сигналов, осуществляющий оцифровку и передачу токов и напряжений через SV (МЭК-61850).\par\vspace{0.2cm}

% DISS-005-B0770
Далее в работе будут подробно рассмотрены вопросы применения КР для работы РНМ.\par\vspace{0.2cm}

% DISS-005-B0771
\subsection{Теоретическое исследование и сравнение КР и ТТ}\label{subsec:ch2_theoretical_study}

% DISS-005-B0772
В первой части проводится сравнение поведения КР и ТТ в идеальных условиях. Моделирование проводилось с использованием пакета MATLAB/Simulink, где анализировались сигналы первых гармоник тока. Схема получения исследуемых сигналов представлена на рисунке ниже (Рисунок 2.1Ошибка! Источник ссылки не найден.).\par\vspace{0.2cm}

% DISS-005-B0773 (image54.png, DISS-005-FIG-051, TAB-037)
% DISS-005-B0774 (пустая ячейка таблицы TAB-037)
% DISS-005-B0775 (подпись Рисунка 2.1)
\begin{figure}[htbp]
    \centering
    \includegraphics[width=0.85\textwidth]{image54.png}
    \caption[Схема получения исследуемых сигналов]{Схема получения исследуемых сигналов: $i(t)$~--- мгновенные значения первичного тока, $i_{\text{CT}}(t)$~--- мгновенное значение тока с идеального ТТ, $u_{\text{CR}}(t)$~--- мгновенное значение напряжения на катушке КР с учётом масштабирования, $k$~--- коэффициент трансформации, ДПФ~--- дискретное преобразование Фурье, $h_1$~--- комплексное значение вектора первой гармоники тока}
    \label{fig:ch2_signal_scheme}
\end{figure}

% DISS-005-B0776
Примечание: Коэффициент трансформации (k) задан равным 1 и в дальнейших расчётах не учитывается.\par\vspace{0.2cm}

% DISS-005-B0777

% DISS-005-B0778
В качестве входного сигнала использовался синусоидальный сигнал с апериодической составляющей, определяемой согласно первому закону коммутации о неразрывности тока (Ошибка! Источник ссылки не найден.), формула (2.1):

% DISS-005-B0779 (отступ формулы TAB-038)
% DISS-005-B0780 (формула 2.1, DISS-005-MATH-001)
% DISS-005-B0781 (номер формулы 2.1)
\begin{equation}
\label{eq:kz_current}
i(t) = I_{\text{кз}} \cdot \sqrt{2} \cdot \left(\cos\phi_{\text{КЗ}} \cdot e^{-\frac{t}{\tau_0}} - \cos(\omega t + \phi_{\text{КЗ}})\right)
\end{equation}

% DISS-005-B0782
\noindent где $i(t)$~--- мгновенные значения сигнала тока, $I_{\text{кз}}$~--- амплитуда тока КЗ, $\phi_{\text{КЗ}}$~--- угол начала КЗ, $\tau_0$~--- постоянная времени затухания, $\omega = 2 \cdot \pi \cdot f$~--- угловая частота, $t$~--- время переходного процесса.\par\vspace{0.2cm}

% DISS-005-B0783
Для катушек Роговского сигнал пропорционален производной первичного тока ($i'(t)$), формула (2.2):

% DISS-005-B0784 (отступ формулы TAB-039)
% DISS-005-B0785 (формула 2.2, DISS-005-MATH-008)
% DISS-005-B0786 (номер формулы 2.2)
\begin{equation}
\label{eq:rogowski_derivative}
i'(t) = I_{\text{кз}} \cdot \sqrt{2} \cdot \left(-\frac{\cos\phi_{\text{КЗ}}}{\tau_0} \cdot e^{-\frac{t}{\tau_0}} + \omega \cdot \sin(\omega t + \phi_{\text{КЗ}})\right)
\end{equation}

% DISS-005-B0787

% DISS-005-B0788 (Рисунок 2.2, DISS-005-FIG-052, TAB-040)
% DISS-005-B0789 (подпись Рисунка 2.2)
\begin{figure}[htbp]
    \centering
    \includegraphics[width=0.75\textwidth]{image55.png}
    \caption{Исследуемый исходный сигнал первичного тока}
    \label{fig:ch2_primary_current}
\end{figure}

% DISS-005-B0790

% DISS-005-B0791
Для удобства отображения на графиках сигналы ТТ и КР в одном масштабе осуществлено приведение сигналов, полученное путём деления мгновенных значений сигнала на угловую частоту ($\omega$). Коэффициент изменения амплитуды может быть рассчитан по следующей формуле (2.3):

% DISS-005-B0792 (отступ формулы TAB-041)
% DISS-005-B0793 (формула 2.3, DISS-005-MATH-010)
% DISS-005-B0794 (номер формулы 2.3)
\begin{equation}
\label{eq:k_cr}
k_{\text{CR}} = \frac{1}{\omega}
\end{equation}

% DISS-005-B0795
\noindent где $k_{\text{CR}}$~--- коэффициент масштабирования периодической составляющей КР, $\omega = 2 \cdot \pi \cdot f$~--- угловая частота.\par\vspace{0.2cm}

% DISS-005-B0796
Так же осуществлялась корректировка полученных комплексных значений с КР по углу (поворот вектора на 90 градусов). Это позволило проанализировать влияние апериодической составляющей на расчет первой гармоники методом быстрого преобразования Фурье (БПФ). Исследования показали, что апериодическая составляющая сигнала уменьшается в $\tau_0$ раз, а периодическая составляющая увеличивается в $\omega$ раз, что при стандартных значениях ($f = 50$~Гц и $\tau_0 = 50$~мс) приводит к уменьшению доли апериодической составляющей в сигнале (Рисунок 2.3) [33]..\par\vspace{0.2cm}

% DISS-005-B0797 (Рисунок 2.3, DISS-005-FIG-053, TAB-042)
% DISS-005-B0798 (подпись Рисунка 2.3)
\begin{figure}[htbp]
    \centering
    \includegraphics[width=0.75\textwidth]{image56.png}
    \caption{Схема Исследуемые сигналы с ТТ и КР.}
    \label{fig:ch2_tt_kr_signals}
\end{figure}

% DISS-005-B0799
Амплитуда и угол первой гармоники тока и приведенного сигнала производной тока с КР определялись при помощи быстрого преобразования Фурье (БПФ), так как именно эти сигналы чаще всего используются в современных микропроцессорных устройствах. Как показали исследования, в отличии от действующего значения (RMS), амплитуда первой гармоники значительно меньше подвержена влиянию апериодической составляющей. Это связано с тем, что RMS при своём расчёте (интеграл суммы квадратов мгновенных значений) использует одновременно все гармоники, в том числе и постоянную составляющую; в то время как при вычислении БПФ, апериодическая составляющая растекается по всему спектру, и вносит меньшие искажения в одну конкретную гармоники. Было выбрано прямоугольное окно с шириной, равной 1 периоду, и количеством точек на период, равным 96. Исследования на моделях показали, что использование каких-либо других окон, например, треугольного, приводит к значительному увеличению колебаний и времени переходного процесса при наличии апериодической составляющей в сигнале, поэтому их применение не рекомендуется.\par\vspace{0.2cm}

% DISS-005-B0800
При определённом значении постоянной времени затухания ($\tau_0$), апериодическая составляющая существенно искажает расчёт амплитуды и угла первой гармоники, что особенно заметно на исходном сигнале, который получается с ТТ (см. Рисунок 2.4). Для оценки этого влияния предлагается производить расчёт времени переходного процесса, которое определяется как время от начала сигнала до последнего момента выхода амплитуды и угла за допустимые диапазоны, равные 1\% и 1$^\circ$ для класса точности 5P соответственно.\par\vspace{0.2cm}

% DISS-005-B0801 (Рисунок 2.4а, DISS-005-FIG-054, TAB-043)
% DISS-005-B0802 (Рисунок 2.4б, DISS-005-FIG-055, TAB-043)
% DISS-005-B0803 (метка а)
% DISS-005-B0804 (метка б)
% DISS-005-B0805 (подпись Рисунка 2.4)
\begin{figure}[htbp]
    \centering
    \begin{minipage}{0.48\textwidth}
        \centering
        \includegraphics[width=\textwidth]{image57.png}
        \centerline{\small а)}
    \end{minipage}\hfill
    \begin{minipage}{0.48\textwidth}
        \centering
        \includegraphics[width=\textwidth]{image58.png}
        \centerline{\small б)}
    \end{minipage}
    \caption[Переходный процесс при максимальной доле апериодической составляющей]{Переходный процесс при максимальной доле апериодической составляющей и постоянной времени затухания равной 50~мс: а)~мгновенные значения и амплитуды первой гармоники, б)~угол первой гармоники}
    \label{fig:ch2_transient_max_aperiodic}
\end{figure}

% DISS-005-B0806
Как можем видеть, датчики тока на основе катушек Роговского обладают множеством преимуществ. Однако в отличие от электромагнитных ТТ, они более чувствительны к шумам при малых значениях первичных токов. Поэтому актуальным является рассмотрение усреднённого сигнала, который в будущем на цифровых подстанциях может быть получен путём агрегирования сигналов с разных измерительных каналов. Усреднённым будем называть сигнал, который получается за счёт суммирования в комплексной форме первой гармоники сигналов, полученной с ТТ и с КР с фиксированными коэффициентами равными в сумме 1. Для демонстрации полезных свойств данный сигнал будет рассматриваться в простом варианте за счёт задания равных коэффициентов связи (2.4). В будущем коэффициенты могут быть подобраны более оптимально.\par\vspace{0.2cm}

% DISS-005-B0807 (отступ формулы TAB-044)
% DISS-005-B0808 (формула 2.4, DISS-005-MATH-017)
% DISS-005-B0809 (номер формулы 2.4)
\begin{equation}
\label{eq:i_averaged}
\dot{I}_{\text{усред.}} = 0{,}5 \cdot \dot{I}_{\text{ТТ}} + 0{,}5 \cdot (-j) \cdot \dot{I}_{\text{КР}}
\end{equation}

% DISS-005-B0810
На основе полученных данных проведен анализ времени переходного процесса (Рисунок 2.5), определяемого как интервал от момента возникновения сигнала до достижения допустимых отклонений амплитуды и угла первой гармоники (1\% и 1$^\circ$ для класса точности 5Р). Для ТТ переходный процесс длился 127~мс, в то время как для КР~--- не превышал 20~мс. Применение усреднённого сигнала хоть и улучшало время переходного процесса, но столь не значительно~-- 97~мс.\par\vspace{0.2cm}

% DISS-005-B0811 (Рисунок 2.5, DISS-005-FIG-056, TAB-045)
% DISS-005-B0812 (подпись Рисунка 2.5)
\begin{figure}[htbp]
    \centering
    \includegraphics[width=0.75\textwidth]{image59.png}
    \caption[Отображение границ для определения времени переходного процесса]{Отображение границ для определения времени переходного процесса и начала наступление установившегося режима для: КР; ТТ и усреднённого сигнала. Точка~1~--- первое пересечение граничных значений, точка~2~--- окончание процесса, 3~--- окончание процесса для усреднённого сигнала.}
    \label{fig:ch2_transient_time_limits}
\end{figure}
